\chapter{Industry Introduction}
\label{ch:industry} 

\noindent
The Indian Institute of Science (IISc) was established in 1909. Over the years, it has made its place among the best institutes in the world in terms of education and advanced scientific and technological research. With strong principles guiding IISc, it has found a way to pursue core knowledge, applied research for both social and industrial benefits. It has collaborated with several well-renowned companies to find solutions to critical problems. Professors have taken the initiative to establish their start-ups as well to get their research to society. IISc continues to encourage and empower researchers from diverse disciplines to solve pressing challenges that have direct social impacts. \\

\section{Department of Electrical Engineering}
The department of Electrical Engineering at the Indian Institute of Science offers a vibrant environment for research in Electrical Engineering. Established in 1911, it is one of the first few departments at IISc. The vision of the department is to provide the leadership to enable India’s excellence in the field of Electrical Engineering. The department is committed to advancement of the frontiers of knowledge in Electrical Engineering and to provide the students with a stimulating and rewarding learning experience.\\

\noindent
Research in the area of Systems Science started in late fifties. During the sixties and seventies the department had a strong research group in Control Theory and Adaptive control. This evolved into work on Learning Systems and Pattern Recognition. The department also had a history of research in biomedical engineering and signal processing. \\

\subsection{Control and Robotics (Ctrl) Lab}
\noindent The increasing prominence of robotics and autonomous technologies has heightened the demand for controllers that are secure, dependable, and resilient. This field is inherently interdisciplinary, requiring expertise in mechanical engineering, electrical engineering, and computer science. Despite significant progress, developing control systems for robotic platforms—particularly those that are complex, large-scale, or involve multiple agents—remains a significant challenge. \\

\noindent \textbf{Control and Robotics (Ctrl) Lab} is a dynamic and collaborative space where students, researchers, and faculty come together to explore the exciting intersection of control systems and robotics. Lab is focused on pushing the boundaries of what's possible in robot design, motion planning, and intelligent control, and welcomes anyone with a curious mind and a desire to learn to join us on this journey. Ctrl Lab is headed by Dr. Kiran Kumari and conducts cross-disciplinary research focused on designing and deploying advanced control systems in real-world hardware. The prominent research areas the lab focuses on are:
\begin{itemize}
    \item Event-triggered control
    \item Multi-agent systems and control
    \item Nonlinear systems and control
    \item Control of quadrotors
    \item Sliding mode control 
\end{itemize}

\noindent The mission of Ctrl lab is to bridge theoretical advancements with practical implementations, ensuring stability and precision in robotic systems. 
